\section{Limitations}\label{sec:limitations}

\paragraph{No strength results.}
The most important limitation comes first. No v4 training run has been done at a scale
that could measure playing strength. The end-to-end runs in \cref{sec:pipeline-eval}
check that the machinery works; they say nothing about learning efficiency. Every
difference from \katago{} marked $\sim$ or ``new'' in \cref{tab:overview-compare} could
help, do nothing, or hurt. In particular, we make no claim that \gofer{} v4 learns faster
than \gofer{} v3, or that it is comparable to \katago{}.

\paragraph{Untested design choices.}
Four choices are adopted on argument alone:
\begin{itemize}
  \item keeping fast-search rows for the value, ownership and score heads
        (\cref{sec:learner-masking}), which carries a known risk of correlation and
        overfitting;
  \item recency-weighted sampling;
  \item per-step EMA in place of snapshot averaging;
  \item warm-start fine-tuning with a cosine schedule each cycle, in place of \katago{}'s
        continuous training on one run.
\end{itemize}
The last choice has a structural cost: when a candidate is rejected, its training steps
are thrown away, where a continuous learner would have kept improving. The v3 design made
this choice, v4 inherits it, and it should be tested against continuous training with
periodic snapshots.

\paragraph{Missing \katago{} techniques.}
v4 still lacks the following:
\begin{itemize}
  \item \katago{}'s exact policy target pruning, which subtracts forced playouts using the
        PUCT formula (the engine only drops root children with fewer than two visits,
        \cref{sec:selfplay-search});
  \item disabling root noise on fast searches;
  \item the score-belief distribution (pdf and cdf losses);
  \item the board-width-scaled term in global pooling;
  \item Go-specific input features (liberties, ladders, pass-alive areas);
  \item randomisation of rules, komi and board size;
  \item handicap and branching games;
  \item growing the network during a run.
\end{itemize}
In \citet{wu2019accelerating}'s ablations, removing Go-specific features and optimisations
cost a $1.55\times$ training-time factor, and removing forced playouts with pruning cost
$1.25\times$. v4 has only a simplified form of the latter. The \code{gpool} network does not depend on
board size, but the feature schema, the self-play generator and the ONNX export are all
still $9\times9$ only.

\paragraph{Performance.}
On the CPU we measured, training throughput per step is unchanged from v3
(\cref{sec:eval}). The workload is compute-bound, and augmentation and EMA add a little
overhead. The GPU-oriented features (device residency, mixed precision, compile,
channels-last) have been written but not benchmarked on a GPU. The \code{gpool}-6$\times$64
network is smaller on disk but no faster to evaluate on CPU than legacy 6$\times$64.

\paragraph{Validation loss is only a proxy.}
The learner selects \code{best.pt} by validation loss on held-out games. Lower loss on
self-play data from the current champion does not have to mean stronger play. The arena
gate is the real test, and its power is limited by the number of games a free CPU runner
can afford.

\paragraph{Unverified quantisation.}
The int8 export has passed only the structural checks: it loads and produces outputs of
the right shape. Its effect on playing strength is unknown.

\paragraph{Authorship and review.}
The v4 code and this report were produced mostly by AI coding agents under the direction
of the project owner, with tests as the main check. The report has not been externally
peer-reviewed.
